\documentclass[12pt,a4paper]{elex2017}
\usepackage[left=2.5cm,right=2.5cm,bottom=2.5cm,top=2.5cm]{geometry}
\usepackage[english]{babel}
\usepackage[T1]{fontenc}
\usepackage[utf8]{inputenc}
\usepackage[unicode=true]{hyperref}
\usepackage{graphicx}
\usepackage[
    style=authoryear,
    natbib=true,
    defernumbers=true
]{biblatex}
\addbibresource{elex-ontolex.bib}

\setcounter{secnumdepth}{3}
\let\subparagraph\paragraph
%\let\subparagraph\relax

% numbering of sections, format of the title
\makeatletter
% we use \prefix@<level> only if it is defined
\renewcommand{\@seccntformat}[1]{%
  \ifcsname prefix@#1\endcsname
    \csname prefix@#1\endcsname
  \else
    \csname the#1\endcsname\quad
  \fi}
\newcommand\prefix@section{\thesection. }
\makeatother

\setlength{\parindent}{0cm}
\setlength{\parskip}{11pt plus 1pt minus 2pt}
\setlength{\bibhang}{1cm}
\setlength{\baselineskip}{16pt}

%\pagestyle{empty}

%%%%%%%%%%%%%%%%%%%%%%%%%%%%%%%%%%%%%%%%%%%%%%%%%%%%%%%%%%%%%%%%%%%%%%%%%%%%%%%%
% DOCUMENT BODY
%%%%%%%%%%%%%%%%%%%%%%%%%%%%%%%%%%%%%%%%%%%%%%%%%%%%%%%%%%%%%%%%%%%%%%%%%%%%%%%%

\begin{document}
\mainmatter
\title{The OntoLex-Lemon Model: development and applications}
\titlerunning{The OntoLex-Lemon Model}
\author{\bf John P. McCrae$^1$, Paul Buitelaar$^1$, Philipp Cimiano$^{2}$}
\institute{$^1$Insight Centre for Data Analytics, National University of Ireland
Galway\\ $^2$ Cognitive Interaction Technology Excellence Cluster, Bielefeld
University\\
E-mail: john@mccr.ae, paul.buitelaar@insight-centre.org,
cimiano@cit-ec.uni-bielefeld.de}
\toctitle{The OntoLex-Lemon Model: development and applications}

\maketitle

\begin{abstract}
Each article must include an abstract of 150 to 200 words in Latin Modern Roman
11 pt with interlinear spacing of 12 pt. The heading Abstract should be
centered, font Times 10 bold. This short abstract will also be used for
printing a Booklet of Abstracts containing the abstracts of all papers
    presented at the Conference.

\keywords{provide 3--5 keywords, separated by semi-colons}
\end{abstract}

\section{Heading 1}

% 25,000-50,000 Charcters = 8-15 pages

The lemon Model, first proposed in~\cite{mccrae2012interchanging}, has become the primary
model for the representation of lexical data on the Semantic Web and has been
further developed in the context of the W3C OntoLex Community Group. In this
paper we will present the development and future outlooks for this model as well
as briefly detailing some of the current applications of the model. After the
lemon model was developed in the context of the Monnet project, it was decided
that the further development of this model should take place within a forum as
open as possible, which fortunately coincided with the creation of community
groups for W3C. The community group structure provided mailing lists and wikis
for discussion of the model and eventually led to the publishing of the model as
a W3C Report (\cite{cimiano2016lexicon}) and as files in the W3C namespace. The
discussions within the group covered all aspects of the model, however the issue
of semantics was of particular interest to the group and led to a major
innovation in the introduction of a lexical concept, as a distinct element from
the ontology reference. The formal distinction between these is that an ontology
reference is an entity in an ontology, which the word denotes.  As such, the
(ontological) meaning of  the question “When did Prince die?” could be
understood with the ontology predicate deathDate as a reference, but the general
concept of dying refers to an event rather than to a date. This also further
extends the application domain of OntoLex-Lemon, from formal applications such
as question answering and semantic parsing to the representation of general
machine-readable dictionaries, including WordNet and digitized versions of
existing dictionaries.
Thus, the OntoLex-Lemon model has continued to expand in its use cases and we
consider two particular cases here. Firstly, the OntoLex-Lemon model has been
adopted in a variety of online dictionaries and this has provided a common
interface to these dictionaries. This can be exploited to provide a single
access point across multiple dictionaries that is implemented in the back-end by
means of generic SPARQL (a Semantic Web version of SQL) queries. Secondly, we
look at the application of OntoLex-Lemon in the context of the WordNet
Collaborative Interlingual Index (\cite{bond2016cili}), where the model is being
used to provide a single interlingual identifier for every concept in every
language. The application of ontolex-lemon to these two cases will be discussed
in detail at the eLex conference and in a final paper if our contribution is
accepted.
Finally, we consider the future of the OntoLex-Lemon model, which we intend to
continue to develop and have recently identified four areas for extension:
increasing support for use cases of the model in representing digitized
dictionaries, the use of clear and defined data categories to improve
interoperability, the development of a module for representing complex
morphological patterns and finally an extension to support the representation of
diachronic and etymological information within the lexicon. These developments
should further increase the applicability and value of the model to more users.


\section{Acknowledgement} 

Place all acknowledgements (including those concerning research grants and
funding) in a separate section at the end of the article.

\section{References}

%\nocite{*}
\printbibliography[
    type={book},
    notkeyword={dictionary},
    title={Books}
]
\printbibliography[
    type={incollection},
    title={Book Sections}
]
\printbibliography[
    type={inproceedings},
    title={Paper in conference proceedings}
]
\printbibliography[
    type={article},
    title={Journal Articles}
]
\printbibliography[
    type={techreport},
    title={Technical Reports}
]
\printbibliography[
    type={misc},
    title={Web Sites and Other Resources}
]
\printbibliography[
    type={book},
    keyword={dictionary},
    title={Dictionaries}
]


\medskip
\begin{minipage}[t]{\textwidth}
    \noindent This work is licensed under the Creative Commons Attribution
    ShareAlike 4.0 International License.
    \vspace{-2ex}
    \begin{center}%
        \url{http://creativecommons.org/licenses/by-sa/4.0/}\linebreak
        \includegraphics[width=2.33cm]{cc.png}%
    \end{center}
\end{minipage}

\end{document}
% vim: noai nocin nosi inde=:
